\documentclass{article}
\usepackage[T1]{fontenc}
\usepackage{lmodern}
\usepackage{graphicx}
\usepackage{amsmath,amssymb}
\usepackage[a4paper,margin=2cm]{geometry}
\usepackage[absolute,overlay]{textpos}
\setlength{\TPHorizModule}{1mm}
\setlength{\TPVertModule}{1mm}
\usepackage[indonesian]{babel}
\usepackage{hyperref}
\setlength{\emergencystretch}{2em}
% unit: Halaman 5

\begin{document}

% === Halaman 5 (llm) ===

\begin{itemize}
    \item Pembangkitan kunci putaran sangat sederhana.
    \item Kunci eksternal yang panjangnya 256 bit dibagi ke dalam delapan bagian yang masing-masing panjangnya 32 bit.
    \item Delapan bagian ini yang dinamakan $K_1, K_2, \dots, K_8$.
    \item Contoh: $\text{K} = \text{'abcdefghijklmnopqrstuvwxyz123456'}$ \quad $(32 \times 8 \text{ bit} = 256 \text{ bit})$
\end{itemize}

\vspace{5mm}

\begin{minipage}{0.4\textwidth}
    $\text{maka } K1 = \text{'abcd'}$ \\
    $\quad\quad\quad K2 = \text{'efgh'}$ \\
    $\quad\quad\quad K3 = \text{'ijkl'}$ \\
    $\quad\quad\quad K4 = \text{'mnop'}$
\end{minipage}
\begin{minipage}{0.4\textwidth}
    $K5 = \text{'qrst'}$ \\
    $K6 = \text{'uvwx'}$ \\
    $K7 = \text{'yz12'}$ \\
    $K8 = \text{'3456'}$
\end{minipage}

\end{document}
